% =====================================================================
%  4 -- THE IRON CORE, AND ITS COLLAPSE.
%
%  Source: tier3.md lines 1261-1935 in full and in order --
%          0.4   "The onion, the iron core, and the shell that builds it"
%                (1261-1494): 0.4.1 (1263-1310), 0.4.2 (1312-1381),
%                0.4.3 (1383-1417), 0.4.4 (1419-1459), 0.4.5 (1461-1494);
%          0.5   "Core collapse, derived end to end" (1498-1935):
%                0.5.1 (1504-1541), 0.5.2 (1543-1621), 0.5.3 (1623-1649),
%                0.5.4 (1651-1701), 0.5.5 (1703-1735), 0.5.6 (1737-1793),
%                0.5.7 (1795-1849), 0.5.8 (1851-1878), 0.5.9 (1880-1935).
%  Plus, re-homed from the source's Part 0 self-check block (0.10,
%          lines 2714-2755): questions 6, 7, 8 (2731-2735) ->
%          sec:selfcheck-collapse.
%  The onion ASCII diagram (1268-1288) is redrawn as a TikZ figure
%          (fig:onion); its [typical: ...] annotation is carried
%          verbatim in the caption.
%  Cross-references reworded: none.  Cross-references retargeted only
%          through the label registry ("SS0.6.5" -> \S\ref{sec:missingderivative}
%          etc.); "Tier 0 SSI.2", "Tier 2 SS0.3.1", "Tier 2 SSV.2",
%          "Tier 2 SS0.3.2", "Tier 2 SS0.3.5" stay as plain text.
% =====================================================================

\section{The iron core, and its collapse}
\label{part:collapse}

% =====================================================================
\subsection{The onion, the iron core, and the shell that builds it}
\label{sec:ironcore}

\subsubsection{The structure at the end}
\label{sec:onion}

A massive star at the end of its life has an onion structure: each burning
stage leaves its ash behind and moves outward as a shell.

\begin{figure}[H]
\centering
\begin{tikzpicture}[
  font=\small,
  shell/.style={draw, line width=0.7pt, rounded corners=3pt},
  lab/.style={anchor=west, font=\small, inner sep=1pt}]

% ---- nested shells, outermost first ----
\draw[shell, fill=pointrule!3]  (0,0)      rectangle (9.5,7.5);
\node[lab] at (0.25,7.15) {H envelope \hspace{2.2cm} (if not stripped)};
\draw[shell, fill=pointrule!6]  (0.5,0.5)  rectangle (9.0,6.75);
\node[lab] at (0.75,6.40) {He shell};
\draw[shell, fill=pointrule!9]  (1.0,1.0)  rectangle (8.5,6.0);
\node[lab] at (1.25,5.65) {C/O shell};
\draw[shell, fill=pointrule!12] (1.5,1.5)  rectangle (8.0,5.25);
\node[lab] at (1.75,4.90) {O/Ne/Mg};
\draw[shell, fill=pointrule!16] (2.0,2.0)  rectangle (7.5,4.5);
\node[lab] at (2.25,4.15) {Si/S shell};
\draw[shell, fill=warnrule!10, line width=0.9pt] (3.3,2.3) rectangle (6.2,3.7);
\node[align=center, font=\small] at (4.75,3.0)
  {\textbf{Fe CORE}\\ $\sim$1.2--2.0 \Msun};

% ---- the annotation arrow: silicon shell burning feeds the core ----
\draw[{Latex[length=2.6mm]}-, line width=0.9pt]
  (7.5,4.15) -- (10.6,4.15);
\node[anchor=west, align=left, font=\small] at (10.7,4.15)
  {\textbf{SILICON SHELL BURNING}\\ feeds the core};
\end{tikzpicture}
\caption{The onion structure of a massive star at the end of its life:
each burning stage leaves its ash behind and moves outward as a shell;
silicon shell burning feeds the inert iron core. Redrawn from the source's
ASCII diagram. Fe core mass $\sim$1.2--2.0~\Msun{}
\typical[O'Connor \& Ott 2011 Table~1 spans 1.24--2.08~\Msun; Sukhbold,
Woosley \& Heger 2018 max $\approx 2.0$].}
\label{fig:onion}
\end{figure}

Two features drive everything downstream.

\paragraph{(i) The iron core is inert and grows.} Nothing exothermic happens in
it. It is supported by degenerate electrons and it is being \emph{fed} by the
silicon shell burning just outside it, which converts silicon to iron-peak
material and deposits it onto the core (\citealt{timmes1996}: the iron core
``does not grow by radiative diffusion, but by a series of convective shell
burning episodes''; \citealt{suwa2018}). The core mass is therefore a
monotonically increasing function of time, approaching a limit it cannot
exceed.

\paragraph{(ii) The shell is where the interesting nuclear physics happens, and
it is not in a steady state.} Silicon shell burning is episodic --- the shell
ignites, burns, exhausts, contracts, reignites (TWW96: ``usually there are just
one or two such episodes''; \citealt{heger2001} place ``the most important
period for determining core structure'' in silicon shell burning) --- and its
conditions sweep through the regime box rather than sitting at a point. The
project's per-state, imposed-$(T, \rho)$ framing is a model of \emph{one call
inside one zone inside one timestep} of that process.

% ---------------------------------------------------------------------
\subsubsection{The effective Chandrasekhar mass}
\label{sec:mch}

The core collapses when it exceeds the mass its degenerate electrons can
support. Tier~0 \S I.2 and Tier~2 \S 0.3.1 derive the zero-temperature result;
the version that actually applies has at least two corrections
(\citealt{timmes1996} list finite entropy, GR, Coulomb corrections and the
surface boundary pressure).

The cold, fully relativistic limit (\citealt{chandrasekhar1931,
chandrasekhar1939}; \citealt[eq.~3.3.17]{shapiro1983}):
%
\begin{equation}
\Mch \;=\; 4\pi\left(\frac{K'}{\pi G}\right)^{3/2}\!\!\left(-\xi_1^2\theta'(\xi_1)\right),
\qquad K' = \frac{\hbar c}{4}(3\pi^2)^{1/3}\left(\frac{\Ye}{m_u}\right)^{4/3},
\label{eq:mch-general}
\end{equation}
%
with the $n = 3$ Lane--Emden constant $-\xi_1^2\theta'(\xi_1) = 2.018$, giving
%
\begin{equation}
\boxed{\;\Mch \;=\; 5.83\,\Ye^{2}\,\Msun.\;}
\label{eq:mch}
\end{equation}
%
(the coefficient recomputes to 5.825 with CODATA constants;
\citealt[eq.~1]{timmes1996} and \citealt{heger2001} write $5.83\,\Ye^2$).
Numerically: 1.46~\Msun{} at $\Ye = 0.5$, 1.18 at 0.45, 1.03 at 0.42, 0.71 at
0.35 (an earlier version printed three decimals inconsistent with 5.83 at the
0.1\% level; corrected in the 2026-08-18 audit). \derivedhere{} The derivative
is
%
\begin{equation}
\frac{\dd\Mch}{\dd\Ye} \;=\; 11.66\,\Ye\,\Msun \;=\; 5.83\,\Msun\ \text{per unit } \Ye \text{ at } \Ye = 0.5,
\label{eq:dmchdye}
\end{equation}
%
so \textbf{$\delta\Ye = 10^{-3}$ moves \Mch{} by $5.8 \times 10^{-3}$~\Msun},
and $\delta\Ye = 10^{-2}$ moves it by 0.058~\Msun. Hold onto those two
numbers; \S\ref{sec:missingderivative} needs them.

\paragraph{Correction 1 --- finite entropy raises it.} A hot core supports more
mass, because thermal pressure adds to degeneracy pressure:
%
\begin{equation}
\Mcheff \;\approx\; 5.83\,\Ye^{2}\left[1 + \left(\frac{s_e}{\pi \Ye}\right)^{2}\right] \Msun ,
\label{eq:mcheff}
\end{equation}
%
with $s_e$ the electron entropy per baryon in units of $k$
(\citealt[eq.~6]{timmes1996}, following \citealt{baron1990}; the same expression
is eq.~1 of \citealt{heger2001} and is used by \citealt{suwa2018}). At
$s_e \approx 1$ and $\Ye = 0.45$ the bracket is 1.5 --- a 50\% correction, not
a small one ($s_e$ ranges 0.4--1 across a 15~\Msun{} iron core, TWW96, so 50\%
is an upper-end illustration). So the core mass at collapse is a
\emph{two-parameter} function (\Ye{} and entropy), and the entropy is set by
the burning history too (\citealt{heger2001}: ``the entropy can be just as
important as \Ye'').

\paragraph{Correction 2 --- general relativity, Coulomb corrections and the
surrounding pressure lower it} (\citealt{timmes1996}: SR/GR alone reduce
$5.83\,\Ye^2$ to 1.42~\Msun{} at $\Ye = 0.5$). The real collapse threshold is
not a clean crossing of \Mch; it is a stability question about a core embedded
in a star, with GR corrections, and it is decided by the equation of state's
adiabatic index (\S\ref{sec:gamma43}).

The honest summary: \textbf{$\Mch \propto \Ye^2$ is the correct scaling and the
right way to see why \Ye{} matters, but the pre-supernova core mass in a real
model is a several-parameter outcome of the whole burning history, of which
\Ye{} is the one that nuclear burning most directly controls.} Overselling the
formula is a real risk in this project's framing; \S\ref{sec:missingderivative}
is where that risk gets named as an open item.

% ---------------------------------------------------------------------
\subsubsection{Compactness: the structural predictor}
\label{sec:compactness}

The observable-facing quantity in the modern literature is not the core mass
but the \textbf{compactness parameter} (\citealt{oconnor2011}, ApJ 730, 70):
%
\begin{equation}
\xi_{M} \;=\; \frac{M/\Msun}{R(M_{\mathrm{bary}}=M)/1000\ \mathrm{km}}\bigg|_{t=\text{bounce}},
\qquad \text{usually } M = 2.5\,\Msun .
\label{eq:compactness}
\end{equation}
%
It measures how much mass sits inside a given radius --- i.e.\ how steeply the
density falls outside the core (\citealt{sukhbold2016} evaluate it at collapse
onset instead, with little difference). Its usefulness: it correlates with
whether the neutrino-driven mechanism succeeds (\S\ref{sec:delayed}):
\citet{oconnor2011} find explosion the likely outcome for $\xic \lesssim 0.45$
and black-hole formation above; the two-parameter criterion of
\citet{ertl2016} refines it. High compactness $\to$ a lot of material raining
onto the shock at high rate $\to$ harder to explode $\to$ black hole. Low
compactness $\to$ the accretion rate drops sharply after the core $\to$ the
shock revives $\to$ neutron star and a supernova.

The compactness is set by the mass and entropy structure of the carbon/oxygen
(and silicon) shells at collapse, and it is famously \textbf{non-monotonic in
progenitor mass} --- the migration of the C- and O-burning shells with mass,
and convective-versus-radiative core burning, produce large changes in \xic{}
for small changes in initial mass (\citealt[\S 4.6]{oconnor2011};
\citealt[\S 7.1]{sukhbold2016}; \citealt{sukhbold2018}), and hence in the
predicted outcome. That sensitivity is \emph{why} the composition and
energetics of late shell burning are worth computing accurately, and it is the
strongest available argument for this project's existence that does not route
through \Mch.

It is also, honestly, an argument the project does not currently make with
numbers. Nobody here has computed $\partial\xic/\partial\Ye$ or
$\partial\xic/\partial(\text{network size})$. \S\ref{sec:missingderivative} and
register item \reg{T-01} are that gap.

% ---------------------------------------------------------------------
\subsubsection{What pre-supernova \texorpdfstring{\Ye}{Ye} actually looks like}
\label{sec:presnye}

For orientation \typical[pinned in the 2026-08-18 audit to the central-\Ye{}
time series of the 15 and 25~\Msun{} models of Heger et al.\ 2001
(Tables~1--2, Figs~3--5), the iron-core $\Ye(m)$ profile of Timmes, Woosley \&
Weaver 1996 (0.42 at the centre rising to 0.48 at the core edge), and the MESA
15~\Msun{} $Y_{e,c} \approx 0.43$ of Farmer et al.\ 2016]:

\begin{itemize}
\item \textbf{$\Ye = 0.5$ exactly} only for material that has never had a weak
  interaction and started from equal $N$ and $Z$. Not realistic anywhere in a
  real star: even the initial metallicity (\nuc{N}{14} $\to$ \nuc{Ne}{22}
  through He burning; \citealt{woosley2007}) introduces neutron excess.
\item \textbf{$\Ye \approx 0.85$--0.87} in a surviving hydrogen envelope ---
  hydrogen has $Z/A = 1$, so any H-rich layer sits nowhere near 0.5. The
  figures below are for the CO core and inward, i.e.\ everything the network
  is ever called on; the envelope is outside the regime box entirely and is
  only mentioned so the number is not misread.
\item \textbf{$\Ye \approx 0.498$--0.499} in the He-burnt but otherwise
  unprocessed layers (the \nuc{Ne}{22}-bearing C/O material; 0.498 at
  O-depletion in \citealt{heger2001}).
\item \textbf{$\Ye \approx 0.49$--0.50} in the oxygen shell (0.495 at the
  O-shell epoch, \citealt{heger2001}).
\item \textbf{$\Ye \approx 0.42$--0.49} in the silicon-burning region and the
  iron core, with the central value lowest (0.489 at Si ignition $\to$
  0.44--0.45 at core contraction in \citealt{heger2001}; $0.42 \to 0.48$
  across the iron core in \citealt{timmes1996}).
\item \textbf{$\Ye \approx 0.42$--0.43} at the centre at collapse onset in the
  $\sim$15~\Msun{} models (the 25--40~\Msun{} LMP models of \citealt{heger2001}
  end at 0.445--0.447).
\end{itemize}

The project's regime box specifies \textbf{$0.45 < \Ye < 0.5$}, which brackets
the silicon-burning and iron-core range and stops just above the most extreme
central values. The neutron excess $\eta = 1 - 2\Ye$ is the more natural
variable for the weak-sector physics: $\eta = 0$ at $\Ye = 0.5$, 0.02 at 0.49,
0.06 at 0.47, 0.10 at 0.45. \textbf{The box spans $\eta$ from 0 to 0.1} ---
i.e.\ the entire dynamic range of the quantity runs from zero to ten percent,
which is why relative errors in \Ye{} are a bad figure of merit and absolute
ones are the right ones. \code{CLAUDE.md}'s gates are absolute
($|\dYe| \lesssim 3 \times 10^{-6}$ per step, $5 \times 10^{-3}$ per
trajectory); that is the correct choice and this is the reason.

% ---------------------------------------------------------------------
\subsubsection{Which stars are these, actually?}
\label{sec:whichstars}

Every number in \S\ref{sec:ironcore} is quoted for ``a massive star'', and the
project's one real-track comparison (Tier~2 \S V.2 / \foreignreg{2}{R-07}) uses
a \textbf{single 20~\Msun{} model}. That is one point in a parameter space with
at least four axes, and the axes matter for which observable in
\S\ref{part:observables} even applies.

\paragraph{Mass.} The pre-supernova core mass, compactness and the burning
history are all non-monotonic in initial mass (\S\ref{sec:compactness};
\citealt{ertl2016}; \citealt{sukhbold2018}). A conclusion drawn from a
20~\Msun{} track does not automatically transfer to 12 or 30~\Msun, and the
\emph{shape} of the $(T, \rho)$ locus in the regime box --- which is what
\foreignreg{2}{R-07} measures --- is a function of the progenitor.

\paragraph{Metallicity.} Metallicity sets the initial CNO abundance, hence the
\nuc{Ne}{22} produced during helium burning, hence the \textbf{initial neutron
excess} of the material entering carbon burning (\citealt{woosley2007}). So
$\eta$ at the start of silicon burning is not zero and is
metallicity-dependent. It also sets the line-driven mass-loss rate, hence
whether the hydrogen envelope survives (\citealt{heger2003}: ``the principal
physics connecting the final evolution of a star to its metallicity is its mass
loss'').

\paragraph{Rotation.} Rotation induces chemical mixing (rotational instabilities
feeding composition across boundaries), enlarges cores, and changes the shell
structure and hence the compactness (\citealt{heger2003}; \citealt{heger2000}).
It also produces the angular momentum that some explosion mechanisms need.

\paragraph{Binarity --- the axis that changes the observable.} The majority of
massive stars are in binaries that interact (\citealt{sana2012}: $> 70\%$ of
Galactic O stars will exchange mass with a companion). Interaction strips the
hydrogen envelope, which converts a Type IIP into a Type Ib/c
(\citealt[Table~1]{heger2003}, for the envelope-mass classification). And
\S\ref{sec:ni56}'s Arnett-rule reading of \nuc{Ni}{56} mass \textbf{applies to
the stripped case and not to the H-rich case} --- so the progenitor's binary
history decides which of the two observational readings of the same nuclear
quantity is the right one.

\paragraph{Why this belongs in a study document about a network emulator.} Two
reasons, both about scope discipline.

\begin{enumerate}
\item \textbf{The single-track comparison is a sample of size one.}
  \foreignreg{2}{R-07}'s finding --- that the 20~\Msun{} track sits within
  0.12--0.18~dex of a \emph{line} in the regime box, and that only 35\% / 62\% /
  82\% of the box is within 0.1 / 0.2 / 0.3~dex of a visited point (numbers
  from Tier~2 \S V.2 \derivedthere; not yet a RESULTS.md row) --- is a
  statement about one star. Whether the union over a progenitor grid fills the
  box, or merely traces a family of nearby lines, is the question that decides
  whether the box's uniform sampling measure is defensible. Nobody has run it,
  and it needs no new physics --- just more tracks. Registered as \reg{T-22}.
\item \textbf{``Which observable'' is progenitor-dependent}, so a motivation
  chain that ends at ``the \nuc{Ni}{56} mass sets the peak luminosity'' has
  silently assumed a stripped-envelope progenitor. \S\ref{sec:ni56} flags this;
  \S\ref{sec:whichstars} is where the assumption actually enters.
\end{enumerate}

% =====================================================================
\subsection{Core collapse, derived end to end}
\label{sec:corecollapse}

This is the longest derivation in Part~0, and it is here because the project's
entire justification is that \Ye{} propagates through it. You cannot judge
whether $3 \times 10^{-6}$ per step is a sane requirement without knowing what
the number feeds.

% ---------------------------------------------------------------------
\subsubsection{The instability: \texorpdfstring{$\Gamma_1 < 4/3$}{Gamma1 < 4/3}}
\label{sec:gamma43}

A self-gravitating sphere is dynamically stable against radial perturbations if
its pressure responds to compression strongly enough. The condition, derivable
from a linear perturbation analysis of the hydrostatic equations, is on the
\textbf{adiabatic index}
%
\begin{equation}
\Gamma_1 \;\equiv\; \left(\frac{\partial \ln P}{\partial \ln \rho}\right)_{s} ,
\qquad\text{stability requires}\qquad \bar\Gamma_1 \;>\; \frac{4}{3}.
\label{eq:gamma1}
\end{equation}
%
(the pressure-averaged form is textbook --- \citealt[\S 6.7]{shapiro1983} ---
and the criterion as the trigger of collapse is \citealt{janka2012}.) The 4/3
is not arbitrary. \S\ref{sec:virial} derived it: for a $\gamma = 4/3$ gas the
virial theorem gives $E_{\mathrm{tot}} = 0$, so the star is marginally bound
and an infinitesimal compression neither costs nor releases net energy. Above
4/3 compression costs energy and the star springs back; below, compression
releases energy and the collapse runs away.

An iron core is supported by \textbf{relativistic} degenerate electrons, for
which $\Gamma_1 \to 4/3$ in the ideal limit (\citealt[\S 2.3]{shapiro1983}). A
real core sits a little above that --- partial degeneracy, finite temperature
and the residual electron mass all push $\Gamma_1$ up by a few percent
\provtag{[textbook; magnitude assumed]} --- while general relativity pushes the
\emph{threshold} up by $O(GM/Rc^2)$ (\S\ref{sec:mch}'s Correction 2;
\citealt{chandrasekhar1964}; \citealt{janka2012, janka2017} --- rotation lowers
it slightly). \textbf{Both margins are small, and the point survives: the
pre-collapse iron core is balanced on a knife edge}, not stable-with-margin,
and any physical process that reduces $\Gamma_1$ below the threshold anywhere
in the core tips it.

Two such processes are present, and they are the two runaways; which dominates
depends on the core's entropy (\citealt{muller2016a}: electron capture for the
degenerate low-mass cores, with photodisintegration contributing for the
higher-entropy massive ones).

% ---------------------------------------------------------------------
\subsubsection{The two runaways}
\label{sec:runaways}

\paragraph{(a) Photodisintegration of the iron peak.} Above $T \approx 7$--10~GK
(strongly density-dependent, and \emph{above} the project's regime box;
\citealt{janka2017} places the onset of substantial disintegration where the
central temperature approaches $1~\mathrm{MeV} \approx 10^{10}$~K) the thermal
photon bath dismantles iron-peak nuclei:
%
\begin{equation}
{}^{56}\mathrm{Fe} \;\to\; 13\,{}^{4}\mathrm{He} + 4n
\qquad Q = -124.4\ \mathrm{MeV} \;=\; -2.22\ \mathrm{MeV/nucleon},
\label{eq:fe56photo}
\end{equation}
%
and then, at higher $T$ still, \nuc{He}{4} $\to 2p + 2n$. That second step
costs 7.07~MeV per \emph{helium} nucleon, which is $13 \times 28.30/56 = 6.57$~MeV
per \textbf{iron} nucleon, so the two steps chain to
%
\begin{equation}
2.22 + 6.57 \;=\; 8.79\ \mathrm{MeV/nucleon}
\;=\; \frac{B({}^{56}\mathrm{Fe})}{56},
\label{eq:dissbudget}
\end{equation}
%
the total cost of taking iron all the way to free nucleons --- which is of
course just \nuc{Fe}{56}'s binding energy per nucleon, and the check that the
bookkeeping is right. \derivedhere[from AME2020 \citep{wang2021} binding energies:
$B(\nuc{Fe}{56}) = 492.26$~MeV, $B(\nuc{He}{4}) = 28.30$~MeV]

This is \emph{endothermic}. It consumes thermal energy, so the temperature (and
hence the thermal contribution to pressure) fails to rise as fast as adiabatic
compression would demand. In the $\Gamma_1$ language: energy that should have
gone into pressure goes into breaking nuclei instead, and $\Gamma_1$ drops
below 4/3 (\citealt{janka2017}: ``converts thermal energy to rest-mass energy
\ldots\ causes a reduction of the effective adiabatic index \ldots\ below the
critical value of 4/3'').

\begin{warn}
\textbf{A caution the study notes should carry}, because it is easy to get
wrong. This dismantling is a \emph{collapse} phenomenon, not a silicon-burning
one. Silicon burning at 3--4~GK \textbf{produces} the iron peak, and the full
dissociation to free nucleons that drives $\Gamma_1$ below 4/3 requires
temperatures above the regime box. Tier~0 \S I.2 flags exactly this trap.

\textbf{But do not over-read that into ``the whole box is iron-peak
dominated''.} The top of the box is iron-peak dominated only at its
\emph{high-density} edge. NSE favours the high-entropy, many-particle state as
density falls, so at ($\Tn = 7.9$, $\rho = 10^7$) the equilibrium is already
$\alpha$-rich --- \S\ref{sec:alpharich}'s freeze-out regime --- and this file's
own witness state 646 ($\Tn = 7.78$, $\rho = 2.0 \times 10^7$) measures
\code{he4} = 0.521 against \code{fe56} = 0.378 under fixed MESA
(\S\ref{sec:witness}). The distinction that matters for \S\ref{sec:runaways} is
that neither corner has dissociated to \emph{free nucleons}, which is the step
that costs 8.79~MeV/nucleon.
\end{warn}

\paragraph{(b) Electron capture.} Free protons liberated by (a), and iron-peak
nuclei directly, capture electrons (\citealt{langanke2003a}: both channels
``control the neutronization of the matter''):
%
\begin{equation}
e^- + p \to n + \nu_e, \qquad e^- + (Z,A) \to (Z-1,A) + \nu_e .
\label{eq:ecreactions}
\end{equation}
%
Each capture does two damaging things at once: it \textbf{removes a
pressure-supporting electron} (lowering $P$ at fixed $\rho$), and it emits a
neutrino that --- until trapping sets in (\S\ref{sec:trapping}) ---
\textbf{leaves the star entirely}, carrying away energy and entropy
(\citealt{janka2007}; \citealt{janka2017}). \Ye{} drops, so $\Mch \propto \Ye^2$
drops (\citealt{janka2012}: the inner core ``scales with the instantaneous
Chandrasekhar mass''), so the core that was marginally supportable no longer
is.

The two reinforce: dissociation makes free protons, which are the best EC
targets per particle --- though in aggregate captures on \emph{nuclei} dominate
during infall, because $Y_p \sim 10^{-6}$ (\citealt{langanke2003a};
\citealt{janka2007}: ``what matters is the product of abundance times capture
rate''); EC lowers the pressure, accelerating contraction, raising $T$,
accelerating dissociation. Collapse becomes dynamical.

% ---------------------------------------------------------------------
\subsubsection{The free-fall timescale}
\label{sec:freefall}

Once pressure support is lost, the core falls essentially freely. The
characteristic time is
%
\begin{equation}
\tff \;=\; \sqrt{\frac{3\pi}{32\,G\rho}} ,
\label{eq:tff}
\end{equation}
%
giving \derivedhere[the textbook homogeneous-sphere free-fall time --- Janka
2017 quotes the order-of-magnitude form $\tff \sim 0.004\,\rho_{12}^{-1/2}$~s]

\begin{center}
\begin{tabular}{@{}ll@{}}
\toprule
\textbf{$\rho$ [g\,cm$^{-3}$]} & \textbf{\tff} \\
\midrule
$10^{9}$  & 66 ms \\
$10^{10}$ & 21 ms \\
$10^{12}$ & 2.1 ms \\
\bottomrule
\end{tabular}
\end{center}

\textbf{The whole collapse takes a few hundred milliseconds}
\typical[$\sim 10^2$--$10^3$~ms in simulations from the onset at
$\rho \sim 10^{10}$~g\,cm$^{-3}$, Janka 2017]. Compare with the
silicon-burning episode that set the initial conditions: the ratio is
$\sim 10^5$--$10^7$ (a day to two weeks of silicon burning against 0.1--0.3~s of
collapse \typical). This extreme separation is why the collapse can be treated
as taking the composition as an initial condition --- the nuclear network has
no time to do anything the weak reactions do not do --- and it is why the
accuracy of the \emph{pre-collapse} \Ye{} is what matters.

% ---------------------------------------------------------------------
\subsubsection{Neutrino trapping --- derived, with numbers}
\label{sec:trapping}

Early in the collapse, neutrinos from electron capture stream out freely and
\Ye{} falls fast. At some density they stop escaping, and the \emph{lepton
fraction} \Yl{} freezes (\Ye{} itself keeps falling slightly as $e^- \to \nu_e$
inside the trapped core; \citealt{janka2017}). Deriving where that happens is
worth doing because the trapped lepton fraction is the number that sets the
homologous core mass.

The relevant opacity is \textbf{coherent neutral-current scattering off nuclei}
(\citealt{freedman1974}; \citealt{langanke2003a}; \citealt{janka2017}), whose
cross-section is enhanced by $A^2$ (more precisely $\propto N^2$) because the
whole nucleus scatters in phase when the neutrino wavelength exceeds the
nuclear size:
%
\begin{equation}
\sigma_{\mathrm{coh}} \;\approx\; \frac{\sigma_0}{16}\left(\frac{E_\nu}{m_ec^2}\right)^{2} N^{2},
\qquad \sigma_0 = 1.761\times10^{-44}\ \mathrm{cm}^2 ,
\label{eq:sigmacoh}
\end{equation}
%
(\citealt{janka2017}; the coherent-scattering idea is \citealt{freedman1974};
the full form $A^2[1 - Z/A + (4 \sin^2\theta_W - 1) Z/A]^2$ in place of $N^2$
is textbook --- \citealt[\S 18.4]{shapiro1983} --- not in Janka 2017. An
earlier version of this section wrote $(\sigma_0/4)(E_\nu/m_ec^2)^2 A^2/6$, a
form found in no source and 2.3--2.7$\times$ too large for \nuc{Fe}{56};
corrected in the 2026-08-18 audit.) For $A = 56$, $N = 30$ and $E_\nu = 20$~MeV
this gives $\sigma \approx 1.5 \times 10^{-39}$~cm$^2$ ($1.3 \times 10^{-39}$
with the full factor). The mean free path is $\lambda = 1/(n_A \sigma)$ with
$n_A = \rho N_A/A$, and the neutrinos random-walk out of a core of radius
$R \approx 100$~km in a diffusion time $t_{\mathrm{diff}} \approx 3R^2/(\lambda c)$
($R$ held fixed for the estimate; a 0.5~\Msun{} core has $R \approx 130$~km at
$10^{11}$ and 60~km at $10^{12}$~g\,cm$^{-3}$). Trapping happens when
$t_{\mathrm{diff}}$ exceeds the dynamical time \tff:

\begin{center}
\begin{tabular}{@{}lllll@{}}
\toprule
\textbf{$\rho$ [g\,cm$^{-3}$]} & \textbf{$\lambda$ [cm]} &
\textbf{$t_{\mathrm{diff}}$ [s]} & \textbf{\tff{} [s]} & \textbf{trapped?} \\
\midrule
$10^{11}$ & $6.1 \times 10^{5}$ & 0.016 & 0.0066 & yes (marginally) \\
$10^{12}$ & $6.1 \times 10^{4}$ & 0.16  & 0.0021 & firmly \\
\bottomrule
\end{tabular}
\end{center}

\derivedhere[an order-of-magnitude calculation with fixed $A = 56$,
$E_\nu = 20$~MeV and $R = 100$~km; trapping sets in at a few
$\times\,10^{11}$~g\,cm$^{-3}$ (Janka 2017; Langanke \& Mart\'inez-Pinedo 2003)
and is complete by $\sim 10^{12}$~g\,cm$^{-3}$ (Janka et al.\ 2007; M\"uller
2016), which this reproduces.] \textbf{The $A = 56$ choice biases the answer in
a known direction}: nuclei in the collapsing core are heavier
\typical[$A > 65$ by $\sim 10^{12}$~g\,cm$^{-3}$, Janka et al.\ 2007], and since
$\sigma \propto N^2$ while $n_A \propto \rho/A$, the mean free path goes roughly
as $1/(\rho A)$. Using the real $A$ shortens $\lambda$ and traps
\emph{earlier}, so the densities in the table are upper bounds.

\textbf{Consequences.}

\begin{enumerate}
\item \textbf{The lepton fraction freezes at trapping.} Its value there ---
  $\Yl \approx 0.29$--0.30 with modern capture rates ($Y_{e,c} \approx 0.25$--0.27;
  \citealt{janka2012}; \citealt{muller2016a}); the classic \citealt{bethe1990}
  figure $\Yl \approx 0.36$--0.39 predates the LMP rates (Tier~2 \S 0.3.2) ---
  is the \emph{lepton fraction} \Yl{} that the inner core carries through the
  rest of the collapse (the trapped neutrinos contribute their own pressure,
  and the relevant conserved quantity becomes $\Yl = \Ye + Y_\nu$;
  \citealt{janka2017}). An earlier version of this item put the value ``around
  0.35''; corrected in the 2026-08-18 audit.
\item \textbf{The homologous core mass is set by that trapped value}, through
  the same $M \propto \Yl^2$ relation (\S\ref{sec:homologous};
  \citealt{janka2012}; \citealt{langanke2003a} --- the sources write the
  scaling in \Ye).
\item \textbf{Everything the network does before trapping matters; nothing
  after does.} The chain from silicon burning to explosion runs through a
  single number, evaluated once, at the moment the core stops being able to
  talk to the outside.
\end{enumerate}

That is the sharpest available statement of why this project cares about \Ye.

% ---------------------------------------------------------------------
\subsubsection{The homologous core and the bounce}
\label{sec:homologous}

During collapse the inner part of the core --- where the infall velocity is
subsonic and proportional to radius, $v \propto r$ --- moves
\textbf{homologously}: it collapses self-similarly, maintaining its shape
(\citealt{goldreich1980}; \citealt{yahil1983}; \citealt{janka2017}). Outside a
sonic point the infall is supersonic and cannot communicate; that outer
material simply rains down.

The homologous core's mass is again set by a Chandrasekhar-like relation, now
evaluated at the trapped lepton fraction:
%
\begin{equation}
\Mhom \;\approx\; 5.83\,\Yl^{2}\,\Msun \;\times\;(\text{corrections}),
\label{eq:mhom}
\end{equation}
%
which at $\Yl \approx 0.3$ gives $\approx 0.5$~\Msun{} \derivedhere[Janka 2012,
2017 and M\"uller 2016 give 0.4--0.5~\Msun{} for the shock-formation mass; the
older 0.7--0.8~\Msun{} corresponds to Bethe-1990-era \Yl, and an earlier version
of this section used 0.7~\Msun{} --- corrected in the 2026-08-18 audit].
Collapse continues until the central density reaches nuclear saturation,
$\rho \approx 2.7 \times 10^{14}$~g\,cm$^{-3}$ (\citealt{janka2012, janka2017}),
at which point the nuclear force provides sudden, enormous stiffness:
$\Gamma_1$ jumps well above 4/3 (\citealt{janka2012}).

The inner core overshoots, halts, and \textbf{bounces}. The bounce launches a
shock at the \emph{edge of the homologous core} --- i.e.\ at a mass coordinate
of $\approx 0.4$--0.5~\Msun{} (\citealt{janka2012}) --- inside material that is
still falling in at $\sim 0.1c$ near the sonic point and up to $\sim 0.3c$
further out (\citealt{janka2017}).

\textbf{This is where the whole \Ye{} chain lands.} The shock's birth radius,
and how much overlying material it must traverse, are set by \Mhom, which is
set by \Yl, which is set by \Ye{} at collapse onset, which is set by silicon
burning. $\Mhom \propto \Yl^2$ is the squared link; the
$\Yl \leftarrow Y_{e,\mathrm{onset}}$ link is $O(1)$ and set by infall electron
capture (\citealt{janka2017}: \Ye{} drops from $\sim 0.46$ to $< 0.3$ during
infall).

% ---------------------------------------------------------------------
\subsubsection{Why the prompt shock fails --- the energy budget}
\label{sec:promptshock}

The naive picture (a ``prompt'' explosion: the bounce shock propagates out and
blows the star apart) does not work (\citealt{janka2012, janka2017}), and the
arithmetic showing why is short and worth doing.

The shock must traverse the outer iron core --- of order 0.5--1~\Msun{} of
iron-peak material outside the $\sim$0.4--0.5~\Msun{} homologous core
\typical[Janka 2012, 2017]; in fact it stagnates after traversing only
0.3--0.35~\Msun{} (\citealt{janka2012}) or at an enclosed mass of around
1~\Msun{} (\citealt{janka2017}) --- and it \textbf{dissociates every nucleus it
passes through}, because the post-shock temperature is far above the
dissociation threshold (\citealt{janka2017}: ``roughly 8.8~MeV per nucleon or
$1.7 \times 10^{51}$~erg per 0.1~\Msun''). The cost, from
\S\ref{sec:runaways}, is 8.79~MeV per nucleon, illustrated for 0.5~\Msun:
%
\begin{equation}
E_{\mathrm{diss}} \;=\; \frac{0.5\,\Msun}{m_u}\times 8.79\ \mathrm{MeV}
\;=\; 8.4\times10^{51}\ \mathrm{erg}.
\label{eq:ediss}
\end{equation}
%
\derivedhere{} Compare with the energy available to the shock --- a few
$\times\,10^{51}$~erg \typical[inferred from Janka 2012:
$\sim 1.7 \times 10^{51}$~erg per 0.1~\Msun{} dissociated and stagnation after
0.3--0.35~\Msun; the classic estimate is Bethe 1990] --- and with the observed
explosion energy of a typical core-collapse supernova, $\sim 10^{51}$~erg
(\citealt{burrows2021}; \citealt{woosley2007}). \textbf{The dissociation cost
alone exceeds the shock's entire energy budget by a factor of a few --- the
shock's initial energy is spent after only 0.3--0.35~\Msun{}
(\citealt{janka2012})} (an earlier version said ``by roughly an order of
magnitude''; corrected in the 2026-08-18 audit). On top of that, the
dissociated free nucleons are excellent electron-capture targets, so the
shocked material rapidly deleptonises and emits a burst of neutrinos that
carries away more energy still (the ``neutronisation burst'',
$\sim 2 \times 10^{51}$~erg within $\sim 20$~ms at a peak luminosity near
$4 \times 10^{53}$~erg\,s$^{-1}$; \citealt{janka2017}). Whether the burst is
causal for the stall is disputed between the review authorities:
\citet{janka2012} argues stagnation happens before breakout, so the burst is
not causal; \citet{burrows2021} rank the burst first and dissociation second.

The prompt shock \textbf{stalls}, typically within tens of milliseconds, at a
radius of $\sim$100--200~km (\citealt{janka2007}; \citealt{burrows2021}), and
becomes a standing accretion shock: material continues to fall through it and
pile onto the proto-neutron star.

For context, the energy reservoir that exists but is not being used: the
gravitational binding energy of the newly formed neutron star is
%
\begin{equation}
E_{\mathrm{bind}} \;\approx\; \frac{3}{5}\frac{GM^{2}}{R} \;\approx\; 2.6\times10^{53}\ \mathrm{erg}
\qquad (M = 1.4\,\Msun,\ R = 12\ \mathrm{km}),
\label{eq:ebind}
\end{equation}
%
i.e.\ $\sim 96$~MeV per baryon, or $\sim 10\%$ of the rest mass.
\derivedhere[the same Newtonian estimate is Janka 2017's
$E_b \approx 3.6 \times 10^{53}\,(M/1.5\,\Msun)^2\,(R/10~\mathrm{km})^{-1}$~erg,
and GR estimates give $\sim 2.8$--$3 \times 10^{53}$~erg, \citealt{lattimer2001}] \textbf{The explosion needs $< 1\%$ of that (\citealt{janka2007};
$\approx 0.4\%$ for $10^{51}/2.6 \times 10^{53}$ \derivedhere)} --- though
\citet{burrows2021} caution that this ``1\% tolerance'' framing misleads: over
the $\sim 1$~s during which the explosion energy is set, 4--10\% of the emitted
neutrino energy is absorbed in the gain region. The entire core-collapse
supernova problem is the question of how to transfer that small fraction of
the available energy from the neutron star to the overlying material.

% ---------------------------------------------------------------------
\subsubsection{The delayed neutrino-driven mechanism}
\label{sec:delayed}

The standard answer (\citealt{wilson1985}; \citealt{bethe1985} --- the neutrino
idea goes back to \citealt{colgate1966}; reviews \citealt{janka2007},
\citealt{janka2012}): the proto-neutron star radiates its binding energy as
neutrinos of all flavours over $\sim 10$~s (the SN~1987A signal lasted
$\sim 12$~s, and Bayesian fits to it give a total neutrino energy of a few
$\times\,10^{53}$~erg, \citealt{loredo2002}; tens of seconds to full
transparency --- \citealt{janka2017}; \citealt{janka2007}). A small fraction of
those neutrinos --- a few percent of the luminosity during the accretion phase
(\citealt{janka2012}: ``several percent''; \citealt{burrows2021} quote 4--10\%)
--- is \textbf{reabsorbed} in the semi-transparent region behind the stalled
shock (the ``gain region'', \citealt{bethe1985}), through
%
\begin{equation}
\nu_e + n \to p + e^-, \qquad \bar\nu_e + p \to n + e^+ .
\label{eq:gainreactions}
\end{equation}
%
The absorbed energy heats the gain region, raising its pressure, and if the
heating exceeds the cooling for long enough, the shock is revived and the
explosion succeeds (\citealt{janka2007}; \citealt{muller2016a} for the
$\tau_{\mathrm{adv}} > \tau_{\mathrm{heat}}$ criterion). Modern simulations
show the process is aided substantially by multi-dimensional hydrodynamic
instabilities (convection in the gain region --- \citealt{herant1994};
\citealt{burrows1995} --- and the standing accretion shock instability,
\citealt{blondin2003}) that increase the dwell time of material in the heating
region (\citealt{janka2012}).

\paragraph{Where \Ye{} enters, five separate times.} Worth enumerating because
the project's motivation chain usually names only the first:

\begin{enumerate}
\item \textbf{Pre-collapse core mass} --- $\Mch \propto \Ye^2$ (\S\ref{sec:mch}),
  setting how much iron there is to collapse.
\item \textbf{The homologous core mass and shock birth radius} --- via \Yl{} at
  trapping (\S\ref{sec:homologous}; \citealt{janka2012}: shock formation at
  0.4--0.5~\Msun{} for $\Yl \approx 0.29$--0.30).
\item \textbf{The mass the shock must traverse} --- the difference between the
  iron core mass (1) and the homologous core mass (2) is precisely the material
  whose dissociation drains the shock (\S\ref{sec:promptshock}).
\item \textbf{The accretion rate history and compactness} --- set by the shell
  structure the same burning produced (\S\ref{sec:compactness}), which decides
  whether the delayed mechanism has time to work (\citealt{burrows2021};
  \citealt{ertl2016}; \citealt{boccioli2024} on the Si/Si--O interface).
\item \textbf{The ejecta composition} --- the neutron excess of what actually
  escapes, which is what telescopes measure (\S\ref{part:observables}).
\end{enumerate}

Items 1--4 are why the emulator matters for \emph{whether the star explodes};
item 5 is why it matters for \emph{what comes out}. The project's gates are
written against item 5's proxy (\Ye{} itself), because 1--4 have no assembled
derivative --- \S\ref{sec:missingderivative}.

% ---------------------------------------------------------------------
\subsubsection{The mass cut, and the honest caveat}
\label{sec:masscut}

The \textbf{mass cut} is the mass coordinate separating what falls onto the
neutron star from what is ejected. It is not a free parameter of nature but it
\emph{is} effectively a free parameter of piston/thermal-bomb nucleosynthesis
calculations (\citealt{boccioli2024} call it ``an arbitrary quantity''),
because it emerges from a multi-dimensional simulation nobody runs inside a
nucleosynthesis study; calibrated-engine grids such as \citet{sukhbold2016}
compute it, and \citet{woosley2007} argue it is bracketed by the iron-core
edge and by nucleosynthesis and neutron-star-mass constraints. Material just
inside the cut is the most neutron-rich (deepest, most processed), so the cut's
location strongly affects the ejected iron-peak composition and the
\nuc{Ni}{56} mass (\citealt{woosley2007}; \citealt{the1998} for \nuc{Ti}{44}).

\paragraph{Why this matters for how the project frames its claims.} A chain of
the form

\begin{quote}
better network $\to$ better \Ye{} $\to$ better \Mch{} $\to$ better
explodability $\to$ better \nuc{Ni}{56} prediction
\end{quote}

has, at the ``explodability'' and ``\nuc{Ni}{56}'' links, dependencies on the
explosion mechanism and mass cut that are \emph{larger and less well
constrained} than the composition improvement being sold. The defensible
framing is narrower and still strong:

\begin{result}
A stellar-evolution code currently uses a network too small to get the
pre-collapse \Ye{} and shell structure right, and \emph{knows} it
(\S\ref{sec:costmodel}). An emulator that removes that compromise improves the
\textbf{initial condition} every downstream calculation starts from, whatever
mechanism the downstream calculation then assumes.
\end{result}

That version survives referee contact. The longer chain is a motivation, not a
claim, and the study notes should keep the distinction visible --- which is
exactly what Tier~2 \S 0.3.5 / \foreignreg{2}{R-14} is about.

% ---------------------------------------------------------------------
\subsubsection{What ``explodability'' actually means operationally}
\label{sec:explodability}

\reg{T-01} asks for $\partial(\text{explodability})/\partial\Ye$. Before that
can be a well-posed request you have to know what computes explodability, and
the answer explains why the derivative is hard to get and why it may not be the
right thing to ask for.

\paragraph{What a modern core-collapse simulation is.} The state of the art
couples relativistic hydrodynamics to \textbf{neutrino transport} --- the hard
part, since the neutrinos are the energy carrier (\S\ref{sec:delayed}) and
their transport is a 6-dimensional Boltzmann problem (3 space + 2 angle + 1
energy) per flavour (\citealt{muller2016a}; \citealt{couch2020}). Practical
schemes truncate it: two-moment (M1) closures, ray-by-ray approximations, IDSA,
or full Boltzmann in reduced dimensionality (also leakage and flux-limited
diffusion; \citealt{muller2016a} \S on neutrino transport). A 3D simulation of
one progenitor to $\sim 1$~s post-bounce costs millions of node-hours, i.e.\
$\gtrsim 10^7$ core-hours (\citealt{couch2020}; \citealt{burrows2021}: ``a year
on a supercomputer''). \textbf{You cannot run a grid of those.}

\paragraph{So explodability is computed by proxies}, and there are three
families:

\begingroup\small
\begin{tabularx}{\textwidth}{@{}L{3.4cm} Y L{3.6cm}@{}}
\toprule
\textbf{method} & \textbf{what it does} & \textbf{what it costs} \\
\midrule
\textbf{1D with a calibrated parametrisation} (PUSH, STIR, ``neutrino
light-bulb'' / calibrated inner-boundary engines) &
run 1D with an added, calibrated ingredient --- extra heating (PUSH,
\citealt{perego2015}; light-bulb / calibrated engines, \citealt{ugliano2012},
\citealt{sukhbold2016}) or MLT-turbulence terms fitted to 3D (STIR,
\citealt{couch2020}) --- so that models explode when the calibration set says
they should &
grids of $\sim 200$ progenitors are routine (\citealt{sukhbold2016};
\citealt{couch2020}) [per-model cost of hours: assumed] \\
\textbf{semi-analytic criteria} &
compare a heating timescale with an advection timescale
(\citealt{muller2016a}), or evaluate a critical-luminosity condition
(\citealt{burrows1993}; \citealt{pejcha2015}); the \textbf{compactness} \xic{}
(\S\ref{sec:compactness}; \citealt{oconnor2011}) is the crudest of these &
seconds \\
\textbf{full multi-D} & the real calculation &
$\gtrsim 10^7$ core-hours each (\citealt{couch2020}) \\
\bottomrule
\end{tabularx}
\endgroup

\textbf{The crucial structural fact: explodability is a threshold on a
non-monotonic function.} 1D simulations of most progenitors do \emph{not}
explode without help (except the lowest-mass, ECSN-like cores;
\citealt{janka2012}; \citealt{couch2020}; \citealt{burrows2021}); whether a
given star explodes depends on whether the accretion rate drops fast enough at
the right moment, which depends on the shell structure, which is non-monotonic
in progenitor mass (\citealt{ertl2016}; \citealt{sukhbold2018}). The outcome
map ``initial mass $\to$ explodes or not'' is famously \textbf{islands, not a
boundary} (\citealt{ugliano2012}; \citealt{ertl2016}; \citealt{sukhbold2016}
--- though the island locations disagree between methods,
\citealt{boccioli2024}).

Two consequences for the project's motivation.

\paragraph{(i) $\partial(\text{explodability})/\partial\Ye$ is not a derivative
of a smooth function.} Near a threshold it is a delta function; away from one
it is zero. The scientifically meaningful quantity is not the derivative but
something like \emph{the probability that a $\delta\Ye$ perturbation flips the
outcome}, integrated over a progenitor population --- which requires a grid,
which requires the parametrised 1D methods, which are themselves calibrated.
That is a real research programme, not a literature lookup. \textbf{T-01 should
therefore be split}: the \nuc{Ni}{56} and isotopic-ratio derivatives are
extractable from published yield tables (cheap, do it); the explodability
derivative is not (expensive, and the honest move is to stop claiming it).

\paragraph{(ii) The parametrised-1D grids are exactly where the emulator would
be used.} Those grids consume presupernova model sets that were run to core
collapse hundreds of times (\citealt{sukhbold2016} and \citealt{couch2020} both
use the same $\sim 200$ KEPLER progenitors) --- i.e.\ someone pays the network
cost hundreds of times, in the phase where it dominates
(\S\ref{sec:costmodel}). \textbf{That is the deployment story with the clearest
economics}, and it is better than the vaguer ``improve the initial condition''
framing of \S\ref{sec:masscut}, because it names a workload that exists, is
compute-bound, and is run by people who would notice a 20$\times$ speedup. It
should be in the project's motivation and is not. Registered as \reg{T-23}.

% =====================================================================
\subsection{Self-check}
\label{sec:selfcheck-collapse}

Answer without scrolling up.

\begin{selfcheck}
% source 0.10 Q6
\item State the collapse instability condition and explain why 4/3 appears in
  both it and the virial theorem.
% source 0.10 Q7
\item The bounce shock has $\sim 10^{51}$~erg. Compute the cost of dissociating
  0.5~\Msun{} of iron and say what follows.
% source 0.10 Q8
\item Name the five distinct places \Ye{} enters the collapse-to-explosion
  chain.
\end{selfcheck}
